
\subsection{Grupių prezentacijos}

Grupės prazentacija yra būdas apibrėžti grupę,
pateikiant jos generatorius ir sąryšius kuriuos jie tenkina.
Formaliai tokia grupė yra faktor-grupė laisvosios grupės,
generuojamos duotų generatorių.
Toliau pristatysime kobinatorinės grupių teorijos pagrindus,
reikalingus apibrėžti grupės prezentacija.

Tegul $S = \qty{a,b,c,\cdots}$ yra aibė generatorių.
Literatūroje įprasta aibę $S$ vadinti \emph{abėcėle}, 
o jos elementus - \emph{raidėmis}.
Tegul $S^{-1} = \qty{a^{-1},b^{-1},c^{-1},\cdots}$
yra aibė sudaryta iš \emph{atvikštinių raidžių}
(kol kas šios raidės yra tik aibės $S^{-1}$ elementai be papildomos struktūros).
Aibę $S \cup S^{-1}$ vadinsime \emph{išplėstine abėcėle}.
Baigtines sekas $ f_1 f_2 \cdots f_n$ vadinsime \emph{žodžiais},
čia visi $f_i \in S \cup S^{-1}$.
Tegul $e$ žymi tuščia žodį,
tai yra žodį sudarytą iš nulio raidžių.
Tegul tarp dviejų žodžių yra ekvivalntiškumo sąryšis
$f_1 f_2 \cdots f_n \sim f'_1 f'_2 \cdots f'_m$,
jei iš vieno žodžio galime gauti kitą
pakartotinai įterpdami raidės ir jos anvikštinės raidė poras $xx^{-1}$ ir $x^{-1}x$, $x \in S$.
Visi žodžiai tarp kurių yra ekvivalentiškumo sąryšis sudaro ekivalentiškumo klasę.
\begin{definition}
	Laisvoji grupė $F_S$ generuojama aibės $S$ yra sudaryta iš
	\emph{žodžių} ekvivalentiškumo klasių,
	kartu su daugyba duodama žodžių (klasės atstovų) sujungimo.
\end{definition}

\begin{example}
	Parodykime kad $F_S$ tenkina grupės aksiomas.
	\begin{itemize}
		\item Uždarumas:
			sujungus du žodžius gausime nauja baigtinę raidžių seką, kuri taip pat yra žodis.
		\item Asociatyvumas:
			sekų jungimas yra asociatyvus.
		\item Egzistuoja vienetinis elementas:
			betkurį žodžį sujungus su tuščiu žodžiu $e$, gausime nepakeistą žodį.
		\item Egzistuoja atvikštiniai elementai:
			žodžiui $f_1 f_2 \cdots f_n$ atvikštinis yra $f_{n}^{-1} f_{n-1}^{-1} \cdots f_1^{-1}$,
			kur $f_i^{-1} = x^{-1}$, jei $f_i = x$ ir $f_i^{-1} = x$ jei $f_i=x^{-1}$, $x \in S$.
	\end{itemize}
\end{example}
